\begin{figure*}[t]
\small
\begin{minipage}[t]{0.49\textwidth}
\noindent\textbf{Scala source}

\vspace{0.5em}
\begin{lstlisting}[language=Scala,basicstyle=\ttfamily\codesize,morekeywords={@tailrec,then}]
enum List[T]:
  case Nil()
  case Cons(head: T, tail: List[T])

@tailrec def collect[T, S >: T](
  xs: List[T],
  p: S => Boolean,
  acc: List[{v: T with p(v)}]
): List[{v: T with p(v)}] =
  xs match
    case Nil() => acc
    case Cons(x, xs1) =>
      if p(x)
      then collect(xs1, p, Cons(x, acc))
      else collect(xs1, p, acc)
\end{lstlisting}

\end{minipage}%
\hfill%
\begin{minipage}[t]{0.49\textwidth}
\noindent\textbf{Core calculus encoding}

\vspace{1em}


$\mathit{List}[T] \triangleq \mu X.\, \texttt{Unit} + (T, X)$

\vspace{2.5em}

$\Lambda (T :> \bot <: \top).\; \Lambda (S :> T <: \top).$

\quad$\lambda \mathit{xs}{:}\mathit{List}[T].$

\quad$\lambda p{:}(S \to \texttt{Bool}).$

\quad$\lambda \mathit{acc}{:}\mathit{List}[\{v : T \mid p(v)\}].$

\qquad$\textsf{loop}((\mathit{xs},\, \mathit{acc}))\; \mathit{state}.\;$

\qquad\quad$\textsf{match}\; \mathit{state} \;\textsf{with}\; (\mathit{rem},\, \mathit{cacc}) \Rightarrow$

\qquad\qquad$\textsf{match}\; \mathit{rem} \;\textsf{with}$

\qquad\qquad\quad$\textsf{inl}(\_) \Rightarrow \textsf{inr}(\mathit{cacc})$

\qquad\qquad$\mid\; \textsf{inr}(c) \Rightarrow$

\qquad\qquad\quad$\textsf{match}\; c \;\textsf{with}\; (\mathit{hd},\, \mathit{tl}) \Rightarrow$

\qquad\qquad\qquad$\textsf{if}\; p\;\mathit{hd}$

\qquad\qquad\qquad$\textsf{then}\; \textsf{inl}((\mathit{tl},\, \textsf{inr}((\mathit{hd},\, \mathit{cacc}))))$

\qquad\qquad\qquad$\textsf{else}\; \textsf{inl}((\mathit{tl},\, \mathit{cacc}))$

\end{minipage}

\vspace{0.5em}
\caption{A \texttt{collect} function in Scala (left) and its encoding in the core calculus (right). The tail-recursive loop becomes a $\textsf{loop}$ combinator whose body returns $\textsf{inl}(\mathit{state'})$ to continue or $\textsf{inr}(\mathit{result})$ to break. Lists are equi-recursive: $\textsf{inl}(\texttt{unit})$ represents the empty list and $\textsf{inr}((\mathit{hd}, \mathit{tl}))$ a cons cell.}\label{fig:collect}
\end{figure*}
